\documentclass[a4paper,10pt]{report}\input{../0.Préambule/Préambule_cours_prof}\begin{document}

\newpage
\section*{Probabilités et tableau à double entrée}
\addcontentsline{toc}{section}{Probabilités et tableau à double entrée}

\subsection*{Enoncé}
\addcontentsline{toc}{subsection}{Enoncé}

On tire, deux fois de suite et avec remise, une boule dans une urne contenant une boule bleue et deux boules rouges.

\medskip
\noindent En utilisant un tableau à double entrée, déterminer la probabilité de :
\begin{enumerate}
\item Tirer successivement deux boules rouges,
\item Tirer au moins une boule rouge.
\end{enumerate}

\subsection*{Résolution guidé}
\addcontentsline{toc}{subsection}{Résolution guidé}

On réalise un tableau à double entrée présentant en ligne et en colonne les issues possibles pour chaque tirage :

\begin{center}
\renewcommand{\arraystretch}{2.4}
\begin{tabular}{|C{2.4}|*{3}{C{2.4}|}}
\hline 
\rowcolor{lightgray!20} \diagbox{Urne 1}{Urne 2}& & & \\\hline
\cellcolor{lightgray!20}{} & & & \\\hline
\cellcolor{lightgray!20}{} & & & \\\hline
\cellcolor{lightgray!20}{} & & & \\\hline
\end{tabular} 
\end{center}

\begin{enumerate}
\item On compte le nombre d'issues favorables et le nombre d'issues totales

\medskip \noindent\grandscarreaux{\linewidth}{1.6cm}

On applique la formule de la \textbf{Propriété 4}

\medskip \noindent\grandscarreaux{\linewidth}{1.6cm}
\item On considère l'évènement contraire de \og Tirer au moins une boule rouge \fg{} 

\medskip \noindent\grandscarreaux{\linewidth}{1.6cm}

On compte le nombre d'issues favorables et le nombre d'issues totales

\medskip \noindent\grandscarreaux{\linewidth}{1.6cm}

On conclue avec la \textbf{Propriété 4} et la \textbf{Propriété 3}

\medskip \noindent\grandscarreaux{\linewidth}{1.6cm}
\end{enumerate}

\newpage
\section*{Probabilités et arbres des possibles}
\addcontentsline{toc}{section}{Probabilités et arbres des possibles}

\vspace{-3mm}
\setcounter{cprop}{4}
\begin{prop}
Dans un arbre de probabilités, la probabilité d'une issue est égale au produit des probabilités des branches conduisant à cette issue.
\end{prop}

\subsection*{Enoncé}
\addcontentsline{toc}{subsection}{Enoncé}

On dispose dans une urne de $3$ boules rouges et $7$ boules bleues.
On effectue un premier tirage, puis un second sans replacer la première boule.

\medskip
\noindent En utilisant un arbre des possibles, déterminer la probabilité de :
\begin{enumerate}
\item Tirer successivement deux boules rouges,
\item Tirer au moins une boule rouge.
\end{enumerate}

\subsection*{Résolution guidé}
\addcontentsline{toc}{subsection}{Résolution guidé}

\noindent \begin{minipage}{8cm}
On calcule la probabilités de l'événement \\ R : \og On obtiens une boule rouge lors du premier tirage\fg{} 

\medskip \noindent\grandscarreaux{\linewidth}{1.6cm}

\medskip
On calcule la probabilités de l'événement \\ B : \og On obtiens une boule bleue lors du premier tirage\fg{} 

\medskip \noindent\grandscarreaux{\linewidth}{1.6cm}

\end{minipage}
\hspace{2mm}
\vline
\hspace{2mm}
\begin{minipage}{8.4cm}
On réalise un arbre des possibles en complétant les évènements et les probabilités
\begin{center}
% Racine à Gauche, développement vers la droite
\begin{tikzpicture}[xscale=1,yscale=1,scale=0.8]
% Styles (MODIFIABLES)
\tikzstyle{fleche}=[-,>=latex,thick]
\tikzstyle{noeud}=[fill=white]
\tikzstyle{feuille}=[fill=white]
\tikzstyle{etiquette}=[midway,above]
\tikzstyle{etiquette1}=[midway,below]
% Dimensions (MODIFIABLES)
\def\DistanceInterNiveaux{2}
\def\DistanceInterFeuilles{1}
% Dimensions calculées (NON MODIFIABLES)
\def\NiveauA{(0)*\DistanceInterNiveaux}
\def\NiveauB{(1.6666666666666665)*\DistanceInterNiveaux}
\def\NiveauC{(3)*\DistanceInterNiveaux}
\def\InterFeuilles{(-1)*\DistanceInterFeuilles}
% Noeuds (MODIFIABLES : Styles et Coefficients d'InterFeuilles)
\node[noeud] (R) at ({\NiveauA},{(3.5)*\InterFeuilles}) {};
\node[noeud] (Ra) at ({\NiveauB},{(1.5)*\InterFeuilles}) {\quad \quad \quad \quad \quad \quad \quad \quad};
\node[noeud] (Raa) at ({\NiveauC},{(0.5)*\InterFeuilles}) {};
\node[noeud] (Rab) at ({\NiveauC},{(2.5)*\InterFeuilles}) {};
\node[noeud] (Rb) at ({\NiveauB},{(5.5)*\InterFeuilles}) {\quad \quad \quad \quad \quad \quad \quad \quad};
\node[noeud] (Rba) at ({\NiveauC},{(4.5)*\InterFeuilles}) {};
\node[noeud] (Rbb) at ({\NiveauC},{(6.5)*\InterFeuilles}) {};
% Arcs (MODIFIABLES : Styles)
\draw[fleche] (R)--(Ra);
\draw[fleche] (Ra)--(Raa);
%\draw[fleche] (Raa)--(Raaa);
%\draw[fleche] (Raa)--(Raab);
\draw[fleche] (Ra)--(Rab);
%\draw[fleche] (Rab)--(Raba);
%\draw[fleche] (Rab)--(Rabb);
\draw[fleche] (R)--(Rb);
\draw[fleche] (Rb)--(Rba);
%\draw[fleche] (Rba)--(Rbaa);
%\draw[fleche] (Rba)--(Rbab);
\draw[fleche] (Rb)--(Rbb);
%\draw[fleche] (Rbb)--(Rbba);
%\draw[fleche] (Rbb)--(Rbbb);
\end{tikzpicture}
\end{center}
\end{minipage}

\medskip
On calcule la probabilités manquante de l'arbre :

\medskip \noindent\grandscarreaux{\linewidth}{3.2cm}

\begin{enumerate}
\item On applique la formule de la \textbf{Propriété 5}

\medskip \noindent\grandscarreaux{\linewidth}{1.6cm}
\item On considère l'évènement contraire de \og Tirer au moins une boule rouge \fg{} 

\medskip \noindent\grandscarreaux{\linewidth}{1.6cm}

On conclue avec la \textbf{Propriété 5} et la \textbf{Propriété 3}

\medskip \noindent\grandscarreaux{\linewidth}{1.6cm}
\end{enumerate}
\end{document}